\documentclass[11pt,a4paper]{article}
\usepackage[T1]{fontenc}
\usepackage[utf8]{inputenc}
\usepackage[french]{babel}
\usepackage{geometry}
\usepackage{hyperref}
\usepackage{array}
\usepackage{longtable}
\usepackage{enumitem}

\geometry{margin=2.2cm}
\setlist{nosep}

\title{CSM Robot + Force\\Rapport de groupe}
\author{Equipe CSM\\A completer}
\date{Mars 2026}

\begin{document}
\maketitle
\tableofcontents
\newpage

\section{Contexte et objectifs}

Ce document constitue la base du rendu de groupe du TP CSM Robot + Force.
Il decrit le projet logiciel, la prise en main du capteur, la prise en main du robot xArm6, le couplage robot/capteur, l'asservissement en force sur l'axe Z et les essais realises.

\subsection{Equipe}

\begin{itemize}
    \item Membres: [A completer]
    \item Enseignant / encadrant: [A completer]
    \item Seance(s) et date(s): [A completer]
\end{itemize}

\subsection{Livrables attendus}

\begin{itemize}
    \item description rapide de l'IHM
    \item couplage Robot/Force
    \item asservissement en force
    \item application de test mecanique
    \item captures de code commentees
    \item gestion capteur / liaison serie
    \item gestion bras robotique / reseau
    \item exemple de fichier CSV
\end{itemize}

\section{Architecture du projet}

\subsection{Materiel}

\begin{itemize}
    \item Robot xArm6 (6 axes)
    \item Capteur de force via liaison serie
    \item PC de supervision
\end{itemize}

\subsection{Architecture logicielle}

L'application WinForms regroupe:

\begin{itemize}
    \item la gestion du port serie du capteur
    \item la gestion du reseau xArm
    \item l'IHM de test et de supervision
    \item un mode de guidage manuel a effort nul sur Z
    \item un mode de test mecanique avec export CSV
    \item un asservissement externe approche en force sur Z
    \item la reprise des commandes principales de l'interface xArm d'origine
\end{itemize}

\subsection{Structure du depot}

\begin{longtable}{|p{0.30\linewidth}|p{0.60\linewidth}|}
\hline
\textbf{Chemin} & \textbf{Role} \\
\hline
\texttt{RS232-PC/} & application principale WinForms \\
\hline
\texttt{README.md} & vue d'ensemble du projet \\
\hline
\texttt{TP\_Utilisation\_CSM.md} & guide d'installation, d'utilisation et de test \\
\hline
\texttt{livrables/groupe/} & rendu de groupe versionne \\
\hline
\texttt{livrables/perso/} & espace local ignore par Git \\
\hline
\end{longtable}

\section{Description de l'IHM}

L'application est organisee en cinq onglets:

\begin{itemize}
    \item \texttt{Connexions}
    \item \texttt{Capteur}
    \item \texttt{Robot}
    \item \texttt{xARM Base}
    \item \texttt{Essais}
\end{itemize}

\subsection{Connexions}

Cette vue centralise:

\begin{itemize}
    \item la selection et l'ouverture du port COM
    \item la saisie de l'IP du robot
    \item la connexion xArm
    \item l'activation du mouvement
    \item le journal systeme
\end{itemize}

\subsection{Capteur}

Cet onglet permet:

\begin{itemize}
    \item de lire le log brut du port serie
    \item d'envoyer \texttt{M}, \texttt{C}, \texttt{Q}, \texttt{T}, \texttt{U}
    \item d'afficher la derniere ligne interpretee du type \texttt{Reading: valeur unite}
\end{itemize}

\subsection{Robot et xARM Base}

Les commandes manuelles couvrent:

\begin{itemize}
    \item le jog en repere TOOL
    \item la lecture de la pose et des joints
    \item le retour home, le reset et l'arret d'urgence
    \item les fonctions historiques de l'application xArm d'origine
\end{itemize}

\subsection{Essais}

L'onglet regroupe:

\begin{itemize}
    \item les reglages du timer
    \item les seuils de securite
    \item le chemin de sortie CSV
    \item le guidage manuel a effort nul
    \item le test mecanique
    \item l'asservissement en force
\end{itemize}

\section{Gestion du capteur et liaison serie}

Le service serie ouvre un port en \texttt{115200 / 8N1}, ignore les messages de boot et ne traite comme mesure que les lignes conformes au format suivant:

\begin{quote}
\texttt{Reading: -20.032 Kg}
\end{quote}

La commande d'exploitation courante est \texttt{M}. Les commandes \texttt{T}, \texttt{C}, \texttt{Q} et \texttt{U} restent disponibles pour la tare, la calibration et les tests.

\section{Gestion du robot et liaison reseau}

Le pilotage du robot repose sur le wrapper xArm embarque dans le projet et la DLL native du constructeur.
Le robot cible est un xArm6. L'application n'implemente pas de modele inverse externe; elle utilise les primitives du controleur pour appliquer des deplacements relatifs ou des commandes issues de l'IHM.

\section{Couplage Robot/Force}

La logique d'integration suit la sequence suivante:

\begin{enumerate}
    \item acquisition d'une force via \texttt{M}
    \item lecture de la position robot
    \item calcul d'une correction ou d'un deplacement sur Z
    \item envoi d'une commande \texttt{MoveTool}
    \item journalisation de la mesure dans la session
\end{enumerate}

Ce couplage est exploite dans les modes \texttt{Guidage manuel}, \texttt{Mechanical test} et \texttt{Force control}.

\section{Guidage manuel a effort nul}

Le mode de guidage manuel permet au robot de ceder sur l'axe outil Z quand un operateur pousse ou tire sur le capteur.
La logique retenue est simple et robuste:

\begin{itemize}
    \item tare avant demarrage
    \item consigne d'effort fixe a zero
    \item filtrage exponentiel
    \item zone morte autour de zero
    \item correction relative \texttt{MoveTool} sur Z
\end{itemize}

Ce mode ne constitue pas un freedrive 6 axes. Il s'agit d'un guidage manuel 1 axe, coherent avec le capteur scalaire utilise dans le projet.

\section{Application de test mecanique}

Le test mecanique effectue un balayage sur Z avec un timer unique de 300~ms par defaut.
Chaque cycle:

\begin{itemize}
    \item demande une mesure de force
    \item lit la position Z
    \item stocke l'echantillon
    \item applique un \texttt{Delta Z} fixe
    \item s'arrete sur seuil d'effort, perte de connexion ou sortie de fenetre Z
\end{itemize}

\section{Asservissement en force}

L'asservissement en force est externe au robot et base sur:

\begin{itemize}
    \item un filtrage exponentiel des mesures
    \item une erreur entre mesure et consigne
    \item un calcul PID discret
    \item une saturation du \texttt{Delta Z}
    \item une fenetre de securite sur Z
\end{itemize}

Les gains et parametres sont a ajuster en seance selon le capteur et la direction de l'effort.

\section{Validation et resultats}

\subsection{Tests realises}

\begin{itemize}
    \item compilation et lancement de l'application
    \item test capteur avec emulateur
    \item test capteur reel
    \item connexion robot
    \item jog TOOL
    \item guidage manuel
    \item test mecanique
    \item asservissement en force
\end{itemize}

\subsection{Resultats a inserer}

\begin{itemize}
    \item captures d'ecran de l'IHM
    \item extraits de code commentes
    \item courbes force / position
    \item tableau des parametres de test
    \item difficultes rencontrees et corrections apportees
\end{itemize}

\section{Conclusion}

La base logicielle permet de tester le capteur, de piloter le robot, d'executer des essais de force et de position, puis de preparer un rendu complet.
Les champs marques \texttt{[A completer]} doivent etre finalises par l'equipe avant remise.

\appendix
\section{Annexes a joindre}

Les annexes pourront inclure:

\begin{itemize}
    \item un CSV exemple
    \item des captures commentees
    \item une trace de validation
    \item les fonctions retenues en format texte
\end{itemize}

\end{document}
